%% Согласно ГОСТ Р 7.0.11-2011:
%% 5.3.3 В заключении диссертации излагают итоги выполненного исследования, рекомендации, перспективы дальнейшей разработки темы.
%% 9.2.3 В заключении автореферата диссертации излагают итоги данного исследования, рекомендации и перспективы дальнейшей разработки темы.
\begin{enumerate}
  \item Нелинейное снижение размерности (лапласовы собственные карты, автокодировщики) восстанавливает геометрию среды из популяционной активности гиппокампа при свободном поведении, тогда как линейные методы (PCA) не справляются; качество восстановления растёт с опытом животного и числом зарегистрированных нейронов.

  \item Разработан алгоритм оценки эффективной размерности популяционной активности с коррекцией спектра корреляционных матриц для конечных записей; эффективная размерность снижается во время движения животного. Внутренняя размерность той же активности почти постоянна во времени, значимо ниже эффективной и демонстрирует степенное масштабирование с числом нейронов, что подтверждает нелинейную искривлённость многообразия.

  \item Предложен мультиплексный граф рекуррентности популяции~--- способ агрегации индивидуальных рекуррентных графов нейронов в единое представление, впервые применённый к сырым сигналам кальциевого имиджинга \textit{in~vivo}; из одного и того же графа восстанавливаются как геометрия среды (медианная ошибка декодирования положения около 9~см при стороне арены 44~см), так и динамические режимы коллективной активности, связанные со скоростью движения.

  \item Те же методы анализа применены к искусственным рекуррентным нейронным сетям: установлено, что при равной поведенческой эффективности обучение с подкреплением и обучение с учителем формируют качественно различные вычислительные архитектуры.

  \item Разработан информационно-теоретический фреймворк INTENSE для выявления и количественной оценки нейронной селективности в данных кальциевого имиджинга; показано, что около трети наблюдаемой мультиселективности обусловлено корреляциями поведенческих переменных, а замена взаимной информации на линейную корреляцию приводит к потере более двух третей значимых ассоциаций.

  \item Разработан фреймворк MANGO для безградиентной максимизации активации нейронов импульсных нейронных сетей (оптимизация с разложением тензорного поезда, в 3--15 раз быстрее эталонных методов); обнаружена двухфазная динамика формирования селективности, а последующий анализ инструментами DRIADA и INTENSE выявил двухфазную траекторию в информационной плоскости глубоких слоёв и усиление связи индивидуальной селективности с глубиной сети.

  \item Все разработанные методы объединены в открытый программный пакет DRIADA, интегрирующий анализ нейронной активности на индивидуальном, популяционном и сетевом уровнях; автоматический контроль качества выделенных компонент кальциевого имиджинга (Autoinspect) обеспечил воспроизводимую предобработку записей для всех нейрофизиологических разделов работы.
\end{enumerate}
